\scanpage{8}
\mspage{6}
\margin{\(\circ\!\!-\) \deleted{\unsure{2|3}}}
2.2. La sous-extension \(\overline K_e\) de \(\overline K/K\) est galoisienne, et son groupe de Galois est canoniquement isomorphe au groupe de Galois
\[
\pi_0=\operatorname{Gal}(\bar k/k).\tag*{(2.2.1)}
\]
On obtient ainsi une suite exacte
\[
1\mathrel{\unsure{-|\to}} I\longrightarrow\pi\mathrel{-}\pi_0\mathrel{-}1,\tag*{(2.2.2)}
\]
où
\[
I=\operatorname{Gal}(\overline K/\widetilde K)
\]
est ce qu'on appellé le groupe d'\uline{inertie}. \deleted{\unsure{On}} \deleted{C'est aussi} Si on se rappelle que \(\widetilde K=\) corps des fonctions de \(\widetilde V\), \inserted{et que} \(\overline K\) est une clôture séparable de \(\widetilde K\), on voit donc que \(I\) est le groupe de Galois qui correspond à la situation \deleted{d'un} de l'anneau de valuation discrète \uline{strictement local} \(\widetilde V\). Si \(k\) est sép. clos, on aura \(K=\widetilde K\), \(V=\widetilde V\), \(I=\pi\).

\uline{L'objet du séminaire est} plutôt \uline{l'étude des représentations \(\ell\)-adiques du groupe d'inertie} \(I\). (Ce groupe, ou \unsure{son} image par la représentation \(\ell\)-adique envisagée, \deleted{\unsure{prend alors le|peut alors le|provient de la}} est parfois appelé \uline{groupe de monodromie locale}). Il revient au \(\overline{\mathrm m}\) de dire qu'on s'attache surtout à étudier les repr. \(\ell\)-adiques d'un groupe \(\pi\), lorsque \(k\) est sép. clos. Dans certaines situations cependant, on aura une représentation d'une extension
% archive sheet number: 7
